% Ad Familiares 8.1 (ad-familiares-08-01) --- English translation
% Translated by Claude (Anthropic) from the Latin (Perseus).
% Style guide: STYLE.md.

\ciceroLetterOpener{CAELIVS CICERONI S.}

\ciceroSection{1} Since on parting I promised to write you the fullest possible account of all things going on in the city, I have taken care to engage a man who would follow everything so thoroughly that I am afraid this conscientiousness may strike you as excessive. Still, I know how curious you are, and how welcome it is to anyone abroad to be informed even of the smallest of the goings-on at home. But on one point I beg you not to condemn this service of mine as presumption: that I have delegated this labour to another. Not that it would not be sweetest to me, even busy as I am and ---as you know --- as lazy as can be about writing letters, to give my efforts to the keeping of your memory; but the very volume which I have sent you, as I think, makes my excuses easily enough. I do not know whose leisure it would be, not merely to write all this out, but even to notice it at all: it is all in there --- decrees of the Senate, edicts, stories, rumours. If by chance this specimen is less to your taste, let me know, so that I do not put you to trouble at my expense.

\ciceroSection{2} If anything more important is transacted in the commonwealth than those hirelings can follow conveniently --- both how it was carried, and what general opinion came after it, and what hope there is of the result --- I shall write you a careful account. As things stand, there is no great expectation of anything. Those rumours about the assemblies of the men beyond the Po were hot as far as Cumae, but when I reached Rome I heard not even the faintest whisper of the business. Besides, Marcellus, by not having yet brought forward anything about the succession to the Gallic provinces --- and by having put off the motion, as he told me himself, to the Kalends of June --- has thoroughly reproduced the talk current about him when we were in Rome.

\ciceroSection{3} If you have met Pompey, as you wanted, write me carefully what your impression of him was, what he said to you, and what disposition he displayed --- for it is his way to feel one thing and say another, and he is not so clever as not to let his wish show through.

\ciceroSection{4} As for Caesar, frequent and not very pretty reports come about him, but only as whispers. One man says he has lost his cavalry --- which, I think, certainly happened; another that the Seventh Legion has taken a beating; that the man himself is being besieged among the Bellovaci, cut off from the rest of the army. Nothing is yet certain, and even these uncertain reports are not bandied about openly, but recounted as an open secret among the few you know about --- but Domitius, when he tells them, puts his hand to his mouth. On the 24th of May the loafers in front of the rostra --- may it fall on their own heads --- spread the word that you had died. Through the whole city and the whole forum the loudest rumour ran that you had been killed on the road by Q. Pompeius. I, who knew that Q. Pompeius was running a bathing-business at Baulis, and so starved at that, that I felt sorry for him, was not moved; and prayed that with this lie, if any dangers were hanging over you, we might be quit of them. Your Plancus is at Ravenna; though presented with a fat gift by Caesar, he is neither rich nor well set up. Your books on politics are the rage everywhere.
